% SPDX-FileCopyrightText: Copyright (c) 2026 NVIDIA CORPORATION & AFFILIATES. All rights reserved.
% SPDX-License-Identifier: Apache-2.0
%
% ROLE: how a claim of "better" is decided: splits and hold-out axes, the unit of independence,
%   the noise floor, selection hygiene, the pre-registered headline test and what 12 segments can
%   show, the minimum detectable effect, strata, parity, secondary comparisons, robustness checks,
%   live validation and the staged audits. Details (split tables, noise measurements, the
%   completeness audit, Holm thresholds, SLO constants, lag-patch validation, herding guards, guard
%   metrics, the audit table) live in appendix-protocol.tex (app:protocol, app:audits).
% CLAIMS: protocol only. The primary test was pre-registered before any training
%   (CR/facts/HEADLINE_TEST.json, written at calibration); later amendments (A12, A14-A17) add
%   secondary comparisons and the headline-arm rule, all registered before any test evaluation.
%   The significance arithmetic (minimum attainable p, Holm per N) is our own computation of the
%   exact Wilcoxon null distribution, not a campaign file; it is marked as such.
% EVIDENCE: CR/facts/HEADLINE_TEST.json; CR/cells/{train,val,test}.jsonl; CR/cells/SPLIT_MANIFEST.json;
%   CR/facts/test_freeze.json; CR/facts/{noise,noise_calibration,pilot,gate,lag_patch,remote}.json;
%   CR/CONTRACT.md A1, A2, A4, A5, A11-A17; WT/notes/learned-routing/campaign/PLAN.md
%   (campaign control table); CR/facts/STATE.md audit/fix/gate lines; CR/audits/;
%   WT/benchmarks/learned_routing/learned_routing/goodput.py (guard metrics).
% STATUS: final. The protocol was fixed before the test pass and did not change.

\section{Evaluation protocol}
\label{sec:eval}

% src: CR/facts/HEADLINE_TEST.json written ("2026-10-02 (calibration stage, before any training)"),
%      status; CR/CONTRACT.md A12 (05:50), A14 ("registered BEFORE the test stage runs"), A15, A16,
%      A17 ("Registered before facts/finalists.json and before any test evaluation");
%      CR/facts/pilot.json state.test_cells_touched 0; CR/facts/gate.json selection_hygiene
%      (pilot_records_touching_test 0)
% prov: A12, A14 to A17 (the later additions)
One comparison decides the headline: a pre-registered test over independent workload segments, run
once on the frozen test set; every other analysis is secondary or descriptive. The primary test was
registered at calibration on 2026-10-02, before any training run (\cref{app:repro}).
Later additions, the secondary comparisons and the rule that picks the headline learned arm, were
registered before any test cell was evaluated, and none changed the primary test.

\subsection{Splits and hold-out axes}
\label{sec:splits}

% src: CR/cells/{train,val,test}.jsonl; CR/cells/SPLIT_MANIFEST.json segments_by_split, holdout_axes;
%      CR/cells/val.jsonl (4 cells at N = 6)
% prov: A11.3 (the selection-exposed label)
The test set has 60 cells on 12 segments; \cref{tab:splits,tab:test-axes} in \cref{app:protocol}
give the three splits and the test cells by worker count and hold-out axis. Each test cell is held
out along one primary axis (\cref{sec:freeze}); cells at $\Nw = 6$ are labeled
\emph{selection-exposed} because validation also used that count \citep{cawley2010overfitting}.

\paragraph{Unit of independence.}
% src: CR/facts/HEADLINE_TEST.json test_set.segment_note; CR/cells/test.jsonl (cells per segment:
%      1, 2, 2, 2, 2, 3, 4, 5, 5, 6, 13, 15); CR/CONTRACT.md A4 "Statistics"; CR/literature/LESSONS.md
%      LR-11; WT/notes/learned-routing/README.md (agentx:A1+A2 label)
Cells are not independent: all cells cut from one segment share its arrival and length structure,
whatever their worker count, load level or transform. Every statistic below is therefore computed
over segments. The test has 12, holding between 1 and 15 cells each. The FAST25 conversation cells
share the arrival and length skeleton of Mooncake windows w3 and w5 and are averaged into those two
segments. CRN replicates and duplicated trace copies re-draw or lengthen a segment but add no
independent units. One train cell draws from the plays of two AgentX train segments
(\cref{tab:splits} in \cref{app:protocol}); train segments enter no test statistic.
% prov: A4 (duplicated trace copies)

\subsection{Noise floor}
\label{sec:noise-floor}

% src: CR/facts/noise.json why; CR/literature/LESSONS.md LR-02 [setup]
% prov: A1 (identical re-runs cannot measure noise)
Replay is deterministic for a given policy, workload and seed, so a repeated run reproduces every
per-request row and measures nothing. Every policy therefore runs seeded, and noise is measured over
workload replicates, each of which re-draws the arbitrary parts of the workload and the policy seed
(\cref{sec:aisim}).
% src: CR/facts/noise_calibration.json per_cell_default_goodput_cv (48 non-noise cells, median 0.01376),
%      protocol (8 replicates)
Replicate noise is small: over 8 replicates per cell at calibrated loads, $\piref$'s windowed
goodput has a median coefficient of variation of 1.4\% across the 48 train and validation cells.
% src: CR/facts/pilot.json validation_vs_best_heuristic.llm-d-precise-prefix.M1_tuned (pooled_sd
%      0.01510, segment_se 0.01908), mde.note; CR/facts/gate.json audit_flags_for_stage4[1]
%      (mooncake:w1 +0.102)
Policies differ by segment far more than replicates differ: in the pilot's comparison of M1 with
\policy{llm-d-precise-prefix}, the standard error of the mean over 6 validation segments was 0.019,
against a pooled replicate standard deviation of the paired ratio of 0.015, and M1's lead came almost entirely from one
Mooncake segment (\prelim). Segment heterogeneity, not replicate noise, therefore limits what the
study can resolve; hence the segment-level test and minimum detectable effect below.
\Cref{app:protocol} gives the per-cell ranges, the per-segment pilot numbers and the noise rule the
pilot used.

\subsection{Selection hygiene}
\label{sec:hygiene}

% src: CR/facts/gate.json full_plan.objective_cells_repeats (selection, test), selection_hygiene
%      (val_k3_10_unused_by_selection true; repeats_by_split train 0-7, val 0-2); CR/CONTRACT.md A17.2
The data are split into four disjoint uses, so that selection cannot leak into the result.
\begin{enumerate}
  \item Train cells, replicates $\krep = 0,\dots,7$: the CMA-ES objective (\cref{sec:train-objective}).
  \item Validation cells, $\krep = 0,1,2$: selection within a policy (\cref{sec:selection}).
  \item Validation cells, $\krep = 3,\dots,10$: selection across policies, of $\pibase$ and of the
    headline learned arm $\pilearn$; the pilot's review also used them. No within-policy selection
    reads them.
  \item Test cells, $\krep = 0,1,2$: one pass, for a fixed list of finalists.
\end{enumerate}
% src: CR/facts/test_freeze.json (rule, 149 hashes); CR/facts/test_results.json execution.evaluated_once
%      (149/149); CR/audits/final/completeness.md (30,060 test-cell records, all after finalists.json; F7)
The test split was frozen at calibration (\cref{sec:freeze}), and a completeness audit confirmed
that the test cells were evaluated once (\cref{app:protocol}).

\subsection{Headline test (pre-registered)}
\label{sec:headline}

% src: CR/facts/HEADLINE_TEST.json (binding, per_cell_statistic, per_segment_statistic, primary_test,
%      replicates); CR/CONTRACT.md A2.1, A17.2
The headline compares the selected learned policy $\pilearn$ with the best baseline $\pibase$
(\cref{sec:selection}) on the test set. For each test segment $\seg$,
\begin{equation}
  \label{eq:segstat}
  \segstat_{\seg} = \frac{1}{|\seg|}\sum_{\cell \in \seg}\ \frac{1}{3}\sum_{\krep=0}^{2}
    \clr_{\cell,\krep}(\pilearn, \pibase),
\end{equation}
where $|\seg|$ is the number of test cells in the segment, so each segment weighs the same however
many cells it holds. The primary test is a one-sided exact Wilcoxon signed-rank test over the 12
segment statistics, against the alternative that their median is positive, at level 0.05
\citep{demsar2006}.

It compares one learned policy with one baseline, each chosen on validation before the test, so
no multiplicity correction applies. It does not test the stronger claim that the learned policy
beats \emph{every} baseline; that is a secondary intersection-union comparison
(\cref{sec:secondary}).
% src: CR/literature/LESSONS.md LR-11 action 1 (per-baseline intersection-union test); CR/CONTRACT.md A2.1
% prov: the literature review proposed the every-baseline claim as an intersection of per-baseline
%       tests (LR-11); Amendment A2.1 replaced it with the comparison against the validation-selected
%       best baseline

\paragraph{What 12 segments allow.}
% src: own computation (2026-10-03) of the exact null distribution of the Wilcoxon signed-rank statistic:
%      n = 12: min one-sided p = 1/4096; P(W- <= 17) = 0.0461, P(W- <= 18) > 0.05; n = 4: min p = 1/16;
%      n = 5: min p = 1/32. Consistent with CR/literature/LESSONS.md LR-11 ("At n = 5 its minimum p is 1/32")
With 12 segments the smallest attainable one-sided $p$ is $1/4096$. The test rejects at 0.05
exactly when the segments on which the learned policy loses have ranks, by $|\segstat_{\seg}|$,
summing to at most 17; it can therefore tolerate a few losses if they are small. With 4 segments
the exact test cannot reach 0.05 at all (its minimum $p$ is $1/16$), and with 5 it needs every
segment to favor the same policy. Per-family claims, which rest on 2 to 4 segments, are
therefore descriptive only.

\paragraph{Effect size.}
% src: CR/facts/HEADLINE_TEST.json effect_size; CR/literature/LESSONS.md LR-11 (Bouthillier sec. 4,
%      pdf p.8-10: significant if P - CI_min > 0.5, meaningful if P + CI_max > gamma = 0.75)
We report $\Pr(\pilearn \succ \pibase)$, the share of segments with $\segstat_{\seg} > 0$, with a
95\% bootstrap interval that resamples whole segments together with all their cells. The effect is
called significant if the interval's lower end exceeds 0.5 and meaningful if its upper end exceeds
0.75, following \citet{bouthillier2021variance}.\footnote{Next to it we report the interquartile
  mean and the geometric mean of segment ratios with intervals \citep{agarwal2021precipice}, the
  arithmetic mean of per-cell ratios, a per-cell map of which policy wins, and the gap to the
  per-cell best tuned baseline.}
% src (footnote): CR/facts/HEADLINE_TEST.json effect_size.also_report
% prov: "the arithmetic mean of per-cell ratios" was "the contract's original objective"

\paragraph{Minimum detectable effect.}
% src: CR/facts/HEADLINE_TEST.json parity_and_no_regression.mde; CR/facts/pilot.json
%      mde.learned_vs_val_best_default_arm_heuristic (segment_se 0.019075, mde 0.038150, comparator),
%      mde.learned_vs_default (segment_se 0.043676, mde 0.087352); uses: HEADLINE_TEST n2 parity;
%      p2orch.py drift_matrix (mde 0.038); CR/facts/gate.json full_plan.tiers.C_conditional
The MDE is $\max(1\%,\ 2\,\mathrm{SE})$, where SE is the standard error of the paired segment mean
measured on validation at the pilot. Against \policy{llm-d-precise-prefix} at shipped defaults, the
best untuned heuristic on validation then, it is 0.038 in clipped log-ratio (SE 0.019 over 6
segments); against $\piref$ the same rule gives 0.087. Both are \prelim. The MDE is not a
threshold of the headline test. It is used where a decision needs a size: a difference within it
counts as parity at $\Nw = 2$ (\cref{sec:strata}), it set the drift-check rule
(\cref{sec:drift}), and M2 at ranks 3 and 4 trains only if rank 2 beats M1 on validation by at
least the MDE.

\subsection{Strata and parity}
\label{sec:strata}

% src: CR/facts/HEADLINE_TEST.json stratification; CR/cells/test.jsonl (segments per load mode: open 8,
%      closed 6, lanes 4; transform_extrapolation: 10 cells, 7 segments)
\paragraph{Strata.} Every table and the bootstrap are also given per load mode (8 open-loop, 6
closed-loop and 4 agentic-lane segments), per worker count, for the 10 transform-extrapolation
cells (7 segments) as their own stratum, and per family, descriptively. A significance claim for
an individual worker count uses Holm's correction across the six counts.
% src: own arithmetic (2026-10-03), assuming each per-N test is the same exact one-sided Wilcoxon test:
%      Holm first threshold 0.05/6 = 0.00833; min p with 6 segments 1/64 = 0.0156 > 0.00833.
%      Segments per N from tab:test-axes. Thresholds in app:protocol.
Under it, no unseen worker count, each with 6 test segments, can produce a significant per-count
result whatever the data, so results at unseen counts are descriptive or take the parity and
no-regression form below (our arithmetic, not part of the pre-registration; \cref{app:protocol}).

\paragraph{Parity and no regression.}
% src: CR/facts/HEADLINE_TEST.json parity_and_no_regression (n2, tost_margin); CR/literature/LESSONS.md
%      LR-12 (Calibrate-then-Route sec. V-D, pdf p.4-5: ties at width 3, leads at 4, reconverges at 6)
Two kinds of claim say ``no worse''. For \lc{} at $\theta = -\basis{0}$ against $\piref$, and for
any claim of no regression at an unseen worker count, the protocol uses two one-sided tests with a
margin of $\pm 0.5\%$ normalized goodput. At $\Nw = 2$ a difference within the MDE counts as
parity rather than failure, because informed routing tends to tie simpler policies when there are
few workers to choose from \citep{tumkur2026calibrate}.\footnote{The $\pm 0.5\%$ margin is small
  next to the pilot's segment standard error of 0.019 (\prelim), so before the test we expected most
  no-regression claims to remain unestablished.}
% src (footnote): CR/facts/pilot.json mde (segment_se 0.019075)
In the event, the check held at every unseen worker count against both $\piref$ and
\policy{ramjet}: the lower end of the interval exceeded $-\ln 1.005$ (\cref{sec:res-n}).
% src: CR/facts/test_results.json parity_and_no_regression; CR/report/REPORT.md section 3.2

\subsection{Secondary comparisons}
\label{sec:secondary}

% src: CR/CONTRACT.md A2.1, A2.5, A12, A13 (finalists), A15, A16.3-4, A17.2; CR/facts/HEADLINE_TEST.json
%      effect_size.also_report; CR/literature/LESSONS.md LR-11 (Genet sec. 2 and discussion, pdf p.3-4, p.12)
These were registered before the test pass and are reported next to the headline. Only the
headline is a confirmatory test; the rest are descriptive or secondary.
\begin{itemize}
  \item \textbf{Learned-arm variants}: M1 restricted to its default-router start, that is, selected
    over that restart alone (\cref{sec:restarts}), and M1-noaff, with the session coefficient pinned
    at 0, as test finalists.
    % prov: A12 ("M1-default-init"), A15 (M1-noaff)
  \item \textbf{Reference points}: M0, round-robin, and every heuristic at its shipped defaults.
    % prov: A2.5
  \item \textbf{Oracles}: the best baseline per family and the best tuned baseline per cell, which no
    single policy can match. The gap to the per-cell oracle shows how much a per-workload choice of
    heuristic would add, a comparison \citet{xia2022genet} suggest for learned network policies.
    % prov: A2.1
  \item \textbf{The secondary \ais{}-informed league}: arms that read \ais{}-derived estimates at
    runtime, compared with the best \ais{}-informed baseline and with the best baseline overall.
    Every arm is normalized by the same $\piref$, which does not use \ais{}, so the two leagues share
    one scale and we can state how much the \ais{} estimates add. These arms cannot become
    $\pilearn$ (\cref{sec:res-ais}).
    % prov: A16
  \item \textbf{Every tuned baseline and the runner-up.} The learned arm against each of the 11
    tuned baselines, as an intersection-union test, and against the best baseline of another type,
    \policy{llm-d-precise-prefix}, registered because the validation choice between it and
    \policy{ramjet} was a near tie.
    % prov: LR-11 (intersection-union test)
  \item \textbf{Token-weighted goodput}, registered after the concentration finding: the headline
    pipeline with each good in-window request weighted by its prompt plus output tokens, a check that
    the request-count objective does not hide a loss in served work.
    % prov: A19
% src: CR/facts/finalists.json secondary_comparisons (every_tuned_baseline_IUT, best_baseline_of_another_type
%      llmdpp, runner_up_baseline_registered ablabase); CR/facts/tierBC.json handoff_select_test_i9 (3);
%      CR/CONTRACT.md A19 (context: the phase-2 midpoint re-audit's concentration finding)
\end{itemize}

\subsection{Robustness checks}
\label{sec:robustness}

% src: CR/facts/HEADLINE_TEST.json robustness_reported_not_tested; CR/literature/LESSONS.md LR-14 action 4
%      ("claim only gaps larger than the spread the perturbations induce"; Kendall tau of the ranking)
These checks are reported, not tested, for the finalists on the frozen test cells and replicates.
Each asks whether the headline ranking survives a change that training never saw. Agreement between
two rankings of the finalists is summarized by Kendall's rank correlation $\kendall$. As a reading
rule taken from the literature review rather than from the pre-registration, a headline gap smaller
than the spread these perturbations induce is reported as fragile.
% prov: LR-14 (the reading rule)

\paragraph{SLO re-scoring.}
% src: CR/CONTRACT.md A2.2 (scale sweep, "REPORT states where the headline flips"), A14 (alternatives,
%      constants from train records in facts/sla_transfer_spec.json before test scoring, per-SLA outputs);
%      CR/facts/STATE.md audit(build/goodput-normalization r3) (rescore at all 6 scales verified)
Every evaluation keeps its per-request rows, so the good-request indicator (\cref{eq:good}) can be
recomputed under another SLO without new replays or retraining. The scale sweep multiplies
$(\Imax, \Smax)$ by 0.5, 0.75, 1, 1.5, 2 and 3 and states where the headline ranking flips.
Four alternative definitions of a good request, registered before the test pass:
% prov: A14
\begin{itemize}
  \item the mean-ITL bound alone;
  \item the end-to-end slowdown bound alone;
  \item the mean-ITL bound plus a TTFT bound that grows with prompt length, $b + a\,\ISL_j$, where
    $a$ is three times \ais{}'s uncontended per-token prefill time and $b$ is $\piref$'s median TTFT
    at L2 on train cells;
  \item the mean-ITL bound plus an absolute end-to-end bound, $\piref$'s 90th-percentile end-to-end
    latency at L2 on train cells.
\end{itemize}
For each definition we report every finalist's goodput ratio against $\piref$, the paired
segment delta against $\pibase$, and $\kendall$ against the ranking under the trained SLO. None of
them changes selection or the headline test. Their constants were computed per family from train
records and written down before any test record was scored (\cref{app:protocol}).
% src: CR/facts/sla_transfer_spec.json (records 160; families.mooncake.b_ttft_ms 146.74)

\paragraph{Router-state lag.}
% src: CR/CONTRACT.md A5.1, A11.5, A15; CR/facts/lag_patch.json (summary, semantics, worktree, build,
%      test_stage_usage); CR/literature/LESSONS.md LR-14 (Mitzenmacher pdf p.2)
Replay hands the router perfectly fresh state, and load balancing on stale information can behave
very differently \citep{mitzenmacher1998old}. A replay patch delays three kinds of worker-originated
updates by $\lag$ milliseconds of simulated time: KV cache events, prefill completions and request
completions. A delayed update becomes visible at the first router callback at or after its due time;
the router's own bookkeeping of what it has dispatched stays immediate, as in a live router. At
$\lag = 0$ the patched replay reproduces the unpatched one, and at 50 and \SI{200}{ms} it changes
routing (\cref{app:protocol}). The robustness runs re-evaluated 17 finalists on all test cells at
$\lag \in \{10, 50, 200\}$\,ms, and we report whether the headline ranking holds; M1-noaff ran
alongside M1 to test the hypothesis that session affinity matters more under lag. Training never
sees lag.
% src: CR/facts/robustness.json set, conditions, lag0_identity_lag_build_vs_tuning_build
% prov: A15 (the M1-noaff hypothesis); "The test stage re-evaluated" (the robustness runs); the
%       separate-build detail and the robustness pass's lag-0 identity run moved to app:protocol and
%       app:results
% src: CR/literature/LESSONS.md LR-14 action 2, LR-15 action 3; CR/facts/robustness.json m1_lr14_variants
Because a deterministic argmax over stale state can herd requests onto one worker, two guards
against herding, a positive temperature and a hash-home restriction, ran as variants of M1
(\cref{app:protocol,app:results}).

\paragraph{Timing perturbation.}
% src: CR/facts/HEADLINE_TEST.json robustness_reported_not_tested[2] ("speedup_ratio /
%      decode_speedup_ratio 0.8 and 1.2 (PLAN stage 5)"); WT/lib/mocker/src/common/protocols.rs:549-559
%      (decode_speedup_ratio: decode only; effective decode speedup = speedup_ratio x decode_speedup_ratio)
Two timing scale factors of the mock engines are set to 0.8 and 1.2. The first,
\code{speedup\_ratio}, scales every forward pass; the second, \code{decode\_speedup\_ratio},
scales decode only, so the effective decode speedup is their product. This tests sensitivity of the ranking to a global error in \ais{}'s timing,
not to errors confined to particular batch sizes or context lengths. A timing perturbation also
changes what an uncontended request can achieve, so each perturbed record is scored twice: against
the nominal $\Ezero$ and against the perturbation-consistent
$\Ezero' = \text{prefill}/s + \text{decode}/(s\,d)$ for speedup $s$ and decode speedup $d$. Under the
nominal $\Ezero$ at $s = 0.8$, no zero-reuse AgentX request can be good, a degenerate scoring that is
flagged where it appears. The spread rule and $\kendall$ use $\Ezero'$.
% src: CR/audits/phase2-mid/fix.md "sim-exploitation F2" (rescore_timing, E0' pre-registered, AgentX at
%      s0.8 degenerate under nominal E0); CR/facts/phase2_mid_fix.json preregistration.timing_robustness

\paragraph{Live validation.}
A check on GPUs of the relative result, with its own pre-registered analysis, is described in
\cref{sec:live}; it validates the ranking only and never feeds selection.

\subsection{Audits}
\label{sec:audits}

% src: WT/notes/learned-routing/campaign/PLAN.md "Campaign control" table; CR/audits/ (17 audit
%      reports, 6 fix reports through the pilot's review); CR/facts/STATE.md gate(phase3 judge)
%      ("Independent recompute from raw records ... every pilot number reproduces to 4 decimals");
%      CR/audits/final/refute-headline-1.md (bit-for-bit reproduction of the headline);
%      CR/facts/STATE.md audit(...)/fix(...) lines ("VALID, FIXED, 0 rebutted"; through the pilot 10
%      majors, nine fixed, one resolved by excluding the toolagent trace; phase2-mid 5/5; final 2/2)
% prov: the process details (lenses, static and dynamic audits, grading, fixer, re-audit) and the
%      per-stage major counts moved to app:audits
Each stage of the study ended at a checkpoint that independent auditors reviewed before the next
stage started (\cref{app:audits} describes the process and lists every round in \cref{tab:audits}).
Auditors who recomputed results from the raw records with their own code reproduced every pilot
number to four decimals and the headline bit for bit. Every major finding was judged valid and
resolved, by a fix or, once, by excluding a trace; among them were the objective gaming of
\cref{sec:gaming} and the narrowing of the headline's wording (\cref{sec:res-fairness}).

% src: WT/benchmarks/learned_routing/learned_routing/goodput.py guard_metrics; window_guards.py
%      (added by the phase-2 midpoint fix); CR/literature/LESSONS.md LR-13
% prov: "the midpoint fix" (mid-tuning fix)
Every evaluation record carries guard metrics for worker concentration, session splitting,
starvation and length segregation; after the gaming audit, window-basis versions computed over
exactly the scored requests were added (\cref{sec:gaming,app:protocol}).
